%% kapa: exam=1 round=33 year=2022 period=2 form=A subject=accounting subject_ko=회계학 seq=35 number=75
%% source: Q-Net 감정평가사 기출문제 (cst00309) artlSeq=5209997 file=raw/1008/5209997/2022년도 제33회 감정평가사 1차시험 2교시(감정평가 관계법규, 회계학).pdf
%% conversion: kordoc 4.24.51 -> tools/gen1.py (auto); review=proofread-sonnet; flags=none
\begin{question}{75}
(주)감평의 20×1년 제품 A의 생산·판매와 관련된 자료는 다음과 같다.
\begin{framedtext}
{\small\begin{tabular}{@{}lrl@{}}
단위당 판매가격 & ￦25 & \\
단위당 변동제조원가 & 10 & \\
단위당 변동판매관리비 & 6 & \\
연간 총고정제조간접원가 & 1,500 & (감가상각비 ￦200 포함) \\
연간 총고정판매관리비 & 2,500 & (감가상각비 ￦300 포함) \\
\end{tabular}\par}
\end{framedtext}
(주)감평은 변동원가계산을 채택하고 있으며, 감가상각비를 제외한 모든 수익과 비용은 발생 시점에 현금으로 유입되고 지출된다. 법인세율이 20\%일 때 (주)감평의 세후현금흐름분기점 판매량은?
\begin{choices}
  \item 180단위
  \item 195단위
  \item 360단위
  \item 375단위
  \item 390단위
\end{choices}
\end{question}
\answer{4}
